% Figure 33.2 : le même drain, sans et avec PodDisruptionBudget (mesures du chapitre 33).
\documentclass[tikz,border=4pt]{standalone}
\usepackage{quai}
\begin{document}
\begin{tikzpicture}[x=0.36cm,y=1cm,
    lbl/.style={qlbl, font=\scriptsize},
    rep/.style={draw=#1, fill=#1Bg, line width=0.5pt, rounded corners=1pt}]
  % ---- sans PDB
  \node[font=\titrefig\small, anchor=west] at (-6,4.4) {sans PodDisruptionBudget : drain en 4 s};
  \foreach \k/\y in {1/3.6, 2/3.1, 3/2.6} {
    \node[lbl, anchor=east] at (-0.4,\y) {réplique \k};
    \draw[rep=QVert] (-0.2,\y-0.17) rectangle (0,\y+0.17);
    \draw[rep=QBrique] (0,\y-0.17) rectangle (18,\y+0.17);
    \draw[rep=QVert] (18,\y-0.17) rectangle (30,\y+0.17);
  }
  \node[lbl, text=QBrique] at (9,4.0) {aucune réplique prête : image tirée (13 s), démarrage (5 s)};
  \node[lbl, anchor=west] at (0,2.0) {client : 17 requêtes en échec, pendant environ 17 s};
  % ---- avec PDB
  \node[font=\titrefig\small, anchor=west] at (-6,1.1) {avec minAvailable: 2 : drain en 24 s};
  \foreach \k/\y/\t in {1/0.3/0, 2/-0.2/8, 3/-0.7/16} {
    \node[lbl, anchor=east] at (-0.4,\y) {réplique \k};
    \draw[rep=QVert] (-0.2,\y-0.17) rectangle (\t,\y+0.17);
    \draw[rep=QOcre] (\t,\y-0.17) rectangle (\t+8,\y+0.17);
    \draw[rep=QVert] (\t+8,\y-0.17) rectangle (30,\y+0.17);
  }
  \node[lbl, anchor=west] at (0,-1.3) {client : aucune requête en échec ; jamais moins de 2 répliques prêtes};
  % axe
  \draw[QMuted, line width=0.4pt] (0,-1.8) -- (30,-1.8);
  \foreach \t in {0,5,...,30} { \draw[QMuted] (\t,-1.8) -- (\t,-1.9); \node[lbl, below] at (\t,-1.9) {\t\,s}; }
  % légende
  \draw[rep=QVert] (31,3.5) rectangle (32,3.8); \node[lbl, anchor=west] at (32.2,3.65) {prête};
  \draw[rep=QBrique] (31,3.0) rectangle (32,3.3); \node[lbl, anchor=west] at (32.2,3.15) {absente};
  \draw[rep=QOcre] (31,2.5) rectangle (32,2.8); \node[lbl, anchor=west] at (32.2,2.65) {remplaçante,};
  \node[lbl, anchor=west] at (32.2,2.3) {pas encore prête};
\end{tikzpicture}
\end{document}
